% Chapter 1, Section 2 _Linear Algebra_ Jim Hefferon
%  http://joshua.smcvt.edu/linalg.html
%  2001-Jun-09
\section{线性几何}
\textit{如果你已经见过向量的基本内容，那么
本节是一份可选的复习。
不过，后面的内容会用到这里的材料，
所以如果这对你不是复习，那么它就不是可选的。}

在第一节中，我们费了些功夫才证明
解集只有三种类型\Dash 单点集、空集与
无限集。
但
对于含两个方程两个未知元的方程组，从几何上看这一点很容易。
把每个二元方程画成平面中的一条直线，那么这两条直线
可能交于唯一一点，
平行，或者重合为同一条直线。
%<*diagram:ThreeKindsSolutionSets>
\begin{center}
  \begin{minipage}[b]{1.45in}
    \raisebox{-2pt}[8pt][0pt]{\small \begin{tabular}{@{}l}
      \small \textit{唯一解}
    \end{tabular}}
    \begin{center}
      \includegraphics{gr/mp/ch1.3} \\[.75ex]
      \small $\begin{linsys}{2}
                         3x  &+  &2y  &=  &7   \\
                         x   &-  &y   &=  &-1
                       \end{linsys}$
    \end{center}
  \end{minipage}
  \hspace*{0em}
  \begin{minipage}[b]{1.45in}
    \raisebox{-2pt}[8pt][0pt]{\small \begin{tabular}{@{}l}
      \small \textit{无解}
    \end{tabular}}
    \begin{center}
      \includegraphics{gr/mp/ch1.4} \\[.75ex]
      \small $\begin{linsys}{2}
                         3x  &+  &2y  &=  &7   \\
                         3x  &+  &2y  &=  &4
                       \end{linsys}$
    \end{center}
  \end{minipage}
  \hspace*{0em}
  \begin{minipage}[b]{1.45in}
    \raisebox{-2pt}[8pt][0pt]{\small \begin{tabular}[t]{@{}l}
      \textit{无穷多} \\
      \textit{个解}
    \end{tabular}}
    \begin{center}
      \includegraphics{gr/mp/ch1.5}         \\[.75ex]
      \small $ \begin{linsys}{2}
                         3x  &+  &2y  &=  &7   \\
                         6x  &+  &4y  &=  &14
                       \end{linsys}$
    \end{center}
  \end{minipage}
\end{center}
%</diagram:ThreeKindsSolutionSets>
这些图形并不是证明前一节诸结论的捷径，因为那些结论
适用于含任意多个变量的线性方程组。
但它们确实提供了直观上的认识，是看待那些结论的另一种方式。

本节将建立把我们的结论用几何方式表达所需的工具。
特别地，虽然二维情形是我们所熟悉的，
但要推广到含两个以上未知元的方程组，我们将需要一些高维几何。














\subsectionoptional{空间中的向量}
“高维几何”听起来很新奇。
它确实新奇\Dash 有趣而且令人大开眼界。
但它并不遥远，也并非难以企及。

我们先把一维空间定义为 \( \Re \)。
为看出这个定义是合理的，
想象一维空间
\begin{center}
  \includegraphics{gr/mp/ch1.6}
\end{center}
并选取一点标记为 $0$，另选一点标记为~$1$。
\begin{center}
  \includegraphics{gr/mp/ch1.7}
\end{center}
现在，有了刻度和方向，我们便与 \( \Re \) 建立了对应。
例如，要找与 \( +2.17 \) 相配的点，从 \( 0 \) 出发，朝 \( 1 \) 的方向前进，
走过 \( 2.17 \) 倍的距离即可。

这里把大小与方向结合起来的基本想法，
是推广到更高维的关键。

$\Re^n$ 中由一个大小和一个方向组成的对象
称为\definend{向量\/}\index{vector}
（我们使用与前一节相同的词，因为我们将在下文说明
如何用列向量来描述这样的对象）。
我们可以把向量画成有某个长度并指向某个方向的图形。
\begin{center}
  \includegraphics{gr/mp/ch1.8}
\end{center}
把向量定义为由大小和方向组成时，有一个微妙之处\Dash 下面的这两个
\begin{center}
  \includegraphics{gr/mp/ch1.9}
\end{center}
是相等的，尽管它们从不同的位置出发。
它们相等，是因为它们长度相等且方向相同。
再强调一遍：那些向量不仅相似，它们是相等的。

处于不同位置的东西怎么会相等？
把向量想成表示一个位移
（“vector”一词在拉丁语中的意思是“搬运者”或“旅行者”）。
这两个正方形经历的位移相等，尽管它们从不同的位置出发。
\begin{center}
  \includegraphics{gr/mp/ch1.10}
\end{center}
当我们想强调向量并不固定于某处这一性质时，
我们称它们为\definend{自由}\index{vector!free}向量。
于是，这些自由向量是相等的，
因为每一个都是向右一、向上二的位移。
\begin{center}
  \includegraphics{gr/mp/ch1.12}
\end{center}
更一般地，平面中的向量相同当且仅当它们
第一分量的改变量相同且第二分量的改变量相同：~从
\( (a_1,a_2) \)
延伸到 \( (b_1,b_2) \) 的向量等于从
\( (c_1,c_2) \) 到 \( (d_1,d_2) \) 的向量，当且仅当
\( b_1-a_1=d_1-c_1 \) 且 \( b_2-a_2=d_2-c_2 \)。

说“那个假若从 \( (a_1,a_2) \) 出发
就会延伸到 \( (b_1,b_2) \) 的向量”会很不方便。
我们转而把这个向量描述为
\begin{equation*}
  \colvec{b_1-a_1 \\ b_2-a_2}
\end{equation*}
这样，上面所示的“向右一、向上二”的箭头
就表示成这种形式。
\begin{equation*}
  \colvec[r]{1 \\ 2}
\end{equation*}
我们常把箭头画成从原点出发，并说它处于
\definend{规范位置}\index{vector!canonical position}
（或\definend{自然位置}\index{vector!natural position}
或\definend{标准位置}\index{vector!standard position}）。
当
\begin{equation*}
  \vec{v}=\colvec{v_1 \\ v_2}
\end{equation*}
处于规范位置时，它从原点延伸到终点 $(v_1,v_2)$。

我们通常会说“点
\begin{equation*}
  \colvec[r]{1 \\ 2}\text{”}
\end{equation*}
而不是说那个向量的“规范位置的终点”。
% That is, we shall find it convienent to blur the distinction between a point
% in space and the vector that, if it starts at the origin, ends at that
% point.
因此，我们把下面两者都称为 \( \Re^2 \)。
\begin{equation*}
   \set{(x_1,x_2)\suchthat x_1,x_2\in\Re}
   \qquad
   \set{\colvec{x_1 \\ x_2}\suchthat x_1,x_2\in\Re}
\end{equation*}

在前一节中，我们出于代数上的动机定义了向量与向量运算；
\begin{equation*}
   r\cdot\colvec{v_1 \\ v_2}
   =
   \colvec{rv_1 \\ rv_2}
  \qquad
   \colvec{v_1 \\  v_2}
   +
   \colvec{w_1 \\ w_2}
   =
   \colvec{v_1+w_1 \\ v_2+w_2}
\end{equation*}
现在我们可以从几何上理解这些运算。%
\index{vector!sum}\index{vector!scalar multiple}\index{sum!vector}
\index{scalar multiple!vector}
例如，若 \( \vec{v} \) 表示一个位移，
则 \( 3\vec{v}\, \) 表示沿同一方向但三倍远的位移，
而 \( -1\vec{v}\, \) 表示与 \( \vec{v}\, \) 距离相同但方向相反的位移。
\begin{center}
  \includegraphics{gr/mp/ch1.13}
\end{center}
又，当 \( \vec{v} \) 与 \( \vec{w} \) 表示位移时，
\( \vec{v}+\vec{w}\/ \) 表示这两个位移的合成。
\begin{center}
  \includegraphics{gr/mp/ch1.14}
\end{center}
长箭头就是这个合成位移，其含义如下：想象你正
在轮船的甲板上行走。
假设在一分钟内，
船的运动使船相对于海水产生了位移 $\vec{v}$，
而在同一分钟内，你的行走使你相对于甲板产生了位移 $\vec{w}$。
那么 $\vec{v}+\vec{w}\/$ 就是
你相对于海水的位移。

理解向量和的另一种方式是借助
\definend{平行四边形法则}。\index{parallelogram rule}%
\index{vector!sum}
画出由向量 $\vec{v}$ 与 $\vec{w}$ 构成的平行四边形。
那么和 $\vec{v}+\vec{w}$\index{vector!sum}\index{addition of vectors}
沿对角线延伸到对角的顶点。
\begin{center}
  \includegraphics{gr/mp/ch1.15}
\end{center}

上面的图形展示了向量及向量运算
在 \( \Re^2 \) 中的表现。
我们可以推广到 $\Re^3$，或推广到没有图形可用的更高维空间，
做法是下述显然的推广：~从
\( (a_1,\ldots,a_n) \) 出发而到达 \( (b_1,\ldots,b_n) \) 的
自由向量由这一列表示。
\begin{equation*}
  \colvec{b_1-a_1 \\ \vdotswithin{b_1-a_1} \\ b_n-a_n}
\end{equation*}
若向量有相同的表示，则它们相等。
对于点与规范表示终止于该点的那个向量，我们并不过分讲究二者的区分。
\begin{equation*}
  \Re^n=
  \set{\colvec{v_1 \\ \vdotswithin{v_1} \\ v_n}\suchthat v_1,\ldots,v_n\in\Re}
\end{equation*}
并且，我们按分量做加法与标量乘法。

考虑过点之后，我们接下来考虑直线。\index{line}
在 $\Re^2$ 中，过 \( (1,2) \) 与 \( (3,1) \) 的直线
由这个集合中诸向量（的终点）组成。
% \begin{equation*}
%    \set{ \colvec[r]{1 \\ 2}+t\colvec[r]{2 \\ -1}\suchthat t\in\Re}
% \end{equation*}
% That description expresses this picture.
\begin{center}
  \includegraphics{gr/mp/ch1.16}
\end{center}
在这个描述中，与参数~\( t \) 相伴的向量
\begin{equation*}
  \colvec[r]{2 \\ -1}=\colvec[r]{3 \\ 1}-\colvec[r]{1 \\ 2}
\end{equation*}
就是图中画成整条身躯都落在直线上的那一个\Dash 它是这条直线的
一个\definend{方向向量}\index{vector!direction}%
\index{direction vector}。
注意，这条直线上位于 \( x=1 \) 左边的点
要用 \( t \) 的负值来描述。
% Note also that this description of lines generalizes the familiar
% $y=b+mx$ form for lines in the plane.

在 \( \Re^3 \) 中，
过 \( (1,2,1) \) 与 \( (0,3,2) \) 的直线是这种形式的向量（的终点）组成的集合
\begin{center}
  \includegraphics{gr/mp/ch1.17}
  % \includegraphics{gr/mp/ch1.55}
\end{center}
而更高维空间中的直线也以同样的方式处理。

在 $\Re^3$ 中，
一条直线使用一个参数，因此该直线上的质点
可以沿一维自由地来回移动。
一张平面则涉及两个参数。
例如，过点
\( (1,0,5) \)、\( (2,1,-3) \) 与 \( (-2,4,0.5) \) 的平面
由这个集合中诸向量（的终点）组成。
\begin{equation*}
  \set{ \colvec[r]{1 \\ 0 \\ 5}
         +t\colvec[r]{1 \\ 1 \\ -8}
         +s\colvec[r]{-3 \\ 4 \\ -4.5}
       \suchthat t,s\in\Re      }
\end{equation*}
与各参数相伴的列向量来自这些计算。
\begin{equation*}
  \colvec[r]{1 \\ 1 \\ -8}
  =
  \colvec[r]{2 \\ 1 \\ -3}
  -
  \colvec[r]{1 \\ 0 \\ 5}
  \qquad
  \colvec[r]{-3 \\ 4 \\ -4.5}
  =
  \colvec[r]{-2 \\ 4 \\ 0.5}
  -
  \colvec[r]{1 \\ 0 \\ 5}
\end{equation*}
与直线一样，注意我们用负的 $t$、负的 $s$，或两者兼有，
来描述这个平面中的一些点。

微积分教材常用一个线性方程来描述平面。
% \newsavebox{\jhscratchbox}
% \savebox{\jhscratchbox}{\includegraphics{gr/mp/ch1.18}}
% \newlength{\jhscratchlength}\newlength{\jhscratchheight}
% \settowidth{\jhscratchlength}{\usebox{\jhscratchbox}}
\begin{equation*}
  % \usebox{\jhscratchbox}
  \vcenteredhbox{\includegraphics{gr/mp/ch1.18}}
\end{equation*}
% \begin{equation*}
%   P=\set{\colvec{x \\ y \\ z}\suchthat 2x+3y-z=4}
% \end{equation*}
要从这个描述转到向量描述，
把它看成只含一个方程的线性方程组并参数化：\( x=2-y/2-z/2 \)。
\begin{equation*}
   \vcenteredhbox{\includegraphics{gr/mp/ch1.19}}
\end{equation*}
图中以灰色画出的是与 $y$ 和~$z$ 相伴的向量，它们从原点沿 $x$ 轴平移~$2$ 个单位，
从而整条身躯都落在平面内。
于是这两个向量的和，即以黑色画出的那个，
连同平行四边形的其余部分，整条身躯都在平面内。
% \begin{center}
%   \makebox[\jhscratchlength][l]{\includegraphics{gr/mp/ch1.20}}
% \end{center}
% \begin{equation*}
%   P=\set{\colvec[r]{2 \\ 0 \\ 0}
%          +y\cdot\colvec[r]{-1/2 \\ 1 \\ 0}
%          +z\cdot\colvec[r]{-1/2 \\ 0 \\ 1}\suchthat y,z\in\Re}
% \end{equation*}

推广一下，形如
$\set{\vec{p}+t_1\vec{v}_1+t_2\vec{v}_2+\cdots+t_k\vec{v}_k
             \suchthat t_1,\ldots ,t_k\in\Re}$
（其中 \( \vec{v}_1,\ldots,\vec{v}_k\in\Re^n \) 且 $k\leq n$）的集合是一个
\definend{\( k \) 维线性曲面}\index{linear surface}
（或\definend{\( k \) 平坦集}\index{flat, $k$-flat}）。
例如，在 $\Re^4$ 中
\begin{equation*}
  \set{\colvec[r]{2 \\ \pi \\ 3 \\ -0.5}
       +t\colvec[r]{1 \\ 0 \\ 0 \\ 0}
       \suchthat t\in\Re}
\end{equation*}
是一条直线，
\begin{equation*}
  \set{
       \colvec[r]{0 \\ 0 \\ 0 \\ 0}
       +t\colvec[r]{1 \\ 1 \\ 0 \\ -1}
       +s\colvec[r]{2 \\ 0 \\ 1 \\ 0}
       \suchthat t,s\in\Re}
\end{equation*}
是一张平面，而
\begin{equation*}
  \set{
       \colvec[r]{3 \\ 1 \\ -2 \\ 0.5}
       +r\colvec[r]{0 \\ 0 \\ 0 \\ -1}
       +s\colvec[r]{1 \\ 0 \\ 1 \\ 0}
       +t\colvec[r]{2 \\ 0 \\ 1 \\ 0}
       \suchthat r,s,t\in\Re}
\end{equation*}
是一个三维线性曲面。
同样，直观的理解是：直线允许沿一个方向的运动，
平面允许沿两个方向的组合运动，等等。
当线性曲面的维数比空间的维数小一时，
即当在 $\Re^n$ 中我们有一个
$(n-1)$ 平坦集时，
该曲面称为\definend{超平面}。\index{hyperplane}

对线性曲面的描述在维数方面可能使人产生误解。
例如，这个
\begin{equation*}
  L=\set{
       \colvec[r]{1 \\ 0 \\ -1 \\ -2}
       +t\colvec[r]{1 \\ 1 \\ 0 \\ -1}
       +s\colvec[r]{2 \\ 2 \\ 0 \\ -2}
       \suchthat t,s\in\Re}
\end{equation*}
是一个\definend{退化的}平面，因为它实际上是一条直线，
这是因为
这两个向量互为倍数，我们可以省去其中一个。
\begin{equation*}
  L=\set{
       \colvec[r]{1 \\ 0 \\ -1 \\ -2}
       +r\colvec[r]{1 \\ 1 \\ 0 \\ -1}
       \suchthat r\in\Re}
\end{equation*}
在第二章的“线性无关”一节中，我们将看到向量之间的何种关系
会使它们生成的线性曲面退化。

现在我们可以
用几何语言重述前面得到的结论。
第一，含 \( n \) 个未知元的线性方程组的解集
是 \( \Re^n \) 中的一个线性曲面。
具体地说，它是一个 \( k \) 维线性曲面，其中
\( k \) 是该方程组的某个阶梯形形式中
自由变量的个数。
例如，在单个方程的情形，解集是 $\Re^n$ 中的 $n-1$ 维超平面，其中 $n\geq 1$。
第二，齐次线性方程组的解集
是过原点的线性曲面。
最后，我们可以把任一线性方程组的通解集
看成是其相应的齐次方程组的解集
从原点偏移一个向量而来，
这个向量即任一特解。




